\documentclass{ximera}
\title{Derivatives of Polynomials Worksheet}

\newcommand{\pskip}{\vskip 0.1 in}

\begin{document}
\begin{abstract}
Working with polynomials and their derivatives.
\end{abstract}
\maketitle


\begin{question} \label{QPov332}

Find an equation of the tangent line to the curve
\[
     y = \frac{12}{x^2} - 3x + 4
\]
at the point $(2, 1)$. Show all steps.

\end{question}


\begin{question} \label{Qldf3rmm1}
The function $h=f(t)$, $-1 \leq t \leq 1.8$, expresses the altitude (in thousands of feet) of a balloon in terms of the number of hours past noon. Its graph is shown below.

\begin{onlineOnly}
    \begin{center}
\desmos{8pgzf0orms}{900}{600}
\end{center}
\end{onlineOnly}

\href{https://www.desmos.com/calculator/8pgzf0orms}{151: Height of Balloon 54}



\begin{enumerate}
\item Use the sliders in Lines 2 and 3 above to approximate the balloon's rate of ascent at 12:30pm.  

\item Use the sliders in Lines 2 and 3 above to approximate the time(s) when the balloon is ascending at the rate of $4,000$ ft/hr. 

\item Use the sliders in Lines 2 and 3 above to approximate the time(s) when the balloon's rate of ascent is equal to its average rate of ascent between 11:00am and 1:00pm. 

\item Suppose
\[
      h = f(t) = 5 - 2t - \frac{t^2}{4}+t^3 \, , \, -1 \leq t \leq 1.8,
\]
and find the exact answers to parts (a)-(c).
\end{enumerate}

\end{question}

\begin{question} \label{Qpfp3r398}
The function 
\[
  h = g(t) = h_0 - \frac{1}{2}g t^2 
\]
expresses the height of a rock (in meters) dropped from rest from a height of $h_0$ meters above the lunar surface in terms of the number of seconds since it was released.  Here $g$ is a consant.

\begin{enumerate}
\item Find an expression for the rock's speed (in m/sec) in terms of the number of seconds since it was dropped.

\item Express the speed of the rock when it hits the surface in terms of its average speed during the time it takes to reach the surface.
\end{enumerate}

\end{question}


\begin{question} \label{Qpf3e332}
The function $G=f(s)$, $4\leq s \leq 28$, expresses the number of gallons of gas in your car in terms of the trip odometer reading (measured in miles). It's graph is shown below.

\begin{onlineOnly}
    \begin{center}
\desmos{nnoxloqoro}{900}{600}
\end{center}
\end{onlineOnly}

\href{https://www.desmos.com/calculator/nnoxloqoro}{151: Gas as a Function of Distance 99 2c}


\begin{enumerate}
\item Drap point $P$ above to approximate your (instantaneous) gas mileage when the odometer reading is $10$ miles. Gas mileage is measured in miles/gal and is positive.

\item A dashboard display shows the number of miles you have left to drive assuming your gas mileage remains equal to the current gas mileage for the remainder of your trip. Drag point $P$ to approximate this reading when the odometer reads $15$ miles. {\bf No computations}.

\item Now suppose 
\[
 G  = f(s) = \frac{11}{5} + \frac{1}{5000} \left( s^3  - 50s^2 + 300s \right) \, , \, 4 \leq s \leq 28.
\]

\begin{enumerate}

\item Find the exact answers to parts (a) and (b).

\item Find a function 
\[
   r = w(s) \, , \, 4 \leq s \leq 28.
\]
that expresses the dashboard display for number of miles you have left to drive in terms of the odometer reading. 
\end{enumerate}
\end{enumerate}

\end{question}




\end{document}